%Master File:lectures.tex

\lesson{Sequence Matching}
\vspace{-2cm}
\begin{center}
  \includegraphics[height=10cm]{needlemanWunsch-23}
\end{center}

\keywords{Global matching, dynamic programming, global sequence matching}
%%%%%%%%%%%%%%%%%%%%%%% Next Slide %%%%%%%%%%%%%%%%%%%%%%%
\renewcommand{\Outline}{%
\begin{slide}
\section[1]{Outline}

\begin{minipage}{12cm}
  \begin{enumerate}\squeeze
    \outlineitem{Why Align}{wa}
    \outlineitem{Dynamic Programming}{dp}
    \outlineitem{Global Alignment}{ga}
   \end{enumerate}
\end{minipage}\hfill
\begin{minipage}{10cm}
  \vfill
  \includegraphics[width=10cm]{needlemanWunsch-23}
  \vfill
\end{minipage}
\end{slide}
\addtocounter{outlineitem}{1}
}

\setcounter{outlineitem}{1}
\Outline
\toptarget{firstoutline}
%%%%%%%%%%%%%%%%%%%%%%% Next Slide %%%%%%%%%%%%%%%%%%%%%%%

\begin{slide}
\section{Alignment and Evolution}

\begin{PauseHighLight}
  \begin{itemize}
  \item We want to transfer our knowledge about the function of a
    protein though knowledge of similar proteins in different
    organisms\pause
  \item Most proteins found in an organism are close to those of near
    relatives and sometimes not so near relatives\pause
  \end{itemize}
\end{PauseHighLight}
\begin{center}
  \includegraphics[width=0.9\linewidth]{Histone_Alignment}
\end{center}
\end{slide}

%%%%%%%%%%%%%%%%%%%%%%% Next Slide %%%%%%%%%%%%%%%%%%%%%%%

\begin{slide}
\section[-2]{Practical Alignment}

\begin{PauseHighLight}
  \begin{itemize}
  \item Alignment with known genomes is now the first task after
    sequencing\pause
  \item It is done with both DNA and protein\pause
  \item The mutation rate is sufficiently high that we would expect no
    alignment between most species except if there is a selective
    pressure to conserve the structure\pause
  \item For DNA coding proteins there are a number of ``synonymous
    mutations'' that leave the protein unchanged\pause---DNA is less
    conserved than proteins\pause
  \item In proteins some amino acids behave similarly so there is less
    selection pressure against mutations between them\pause
  \end{itemize}
\end{PauseHighLight}

\end{slide}

%%%%%%%%%%%%%%%%%%%%%%% Next Slide %%%%%%%%%%%%%%%%%%%%%%%

\begin{slide}
\section{Synonymous Mutations}

\begin{minipage}{0.48\linewidth}
  \includegraphics[width=\linewidth]{K12-DH10B-codon-usage}
\end{minipage}\hfil
\begin{minipage}{0.48\linewidth}
  \begin{PauseHighLight}
  \begin{itemize}
  \item Some amino acids are rarer than others\pause
  \item But some codon sequences are slightly more popular than their
    synonymous partners\pause
  \item They may be more robust to mutation\pause
  \item They may be slight biases due to the translation, etc.\pause
  \end{itemize}
\end{PauseHighLight}

\end{minipage}

\end{slide}


%%%%%%%%%%%%%%%%%%%%%%% Next Slide %%%%%%%%%%%%%%%%%%%%%%%

\begin{slide}
\section{Amino Acid Similarity}

\begin{PauseHighLight}
  \begin{itemize}
  \item For amino acids the most important distinguishing property is
    whether they are \emph{hydrophobic} (oily) or \emph{hydrophilic}
    (charged)\pause
  \item There are other properties such as available hydrogen bonds,
    etc.\pause
  \item From many observation we can compute a table giving the log-odds
    of one amino acid being substituted by another\pause
  \item Two commonly used substitution tables (matrices) are PAM and
    BLOSUM\pause
  \item Likely substitutions are positive, unlikely ones negative\pause
  \end{itemize}
\end{PauseHighLight}

\end{slide}

%%%%%%%%%%%%%%%%%%%%%%% Next Slide %%%%%%%%%%%%%%%%%%%%%%%

\begin{slide}
\section[-1.8]{Blosum50}

\begin{center}
  \includegraphics[width=0.98\linewidth]{Blosom50}
\end{center}
\end{slide}

%%%%%%%%%%%%%%%%%%%%%%% Next Slide %%%%%%%%%%%%%%%%%%%%%%%

\begin{slide}
\section{Gaps}

\begin{PauseHighLight}
  \begin{itemize}
  \item It is not infrequent that we find good matches where there has
    been an insertion or deletion of some part of the string\pause
  \item We start by considering a linear cost where the cost of a gap is
    proportional to its length\pause
  \item In biology, a linear cost model is a poor description of
    reality\pause
  \item Where there is an insertion, then it is often quite long (this
    is a consequence of the copy mechanism in cells)\pause
  \item We return to dealing with long gaps in the next lecture\pause
  \end{itemize}
\end{PauseHighLight}

\end{slide}

%%%%%%%%%%%%%%%%%%%%%%% Next Slide %%%%%%%%%%%%%%%%%%%%%%%

\begin{slide}
\section{Inexact Sequence Matching}

\begin{PauseHighLight}
  \begin{itemize}
  \item Given two sequences we want to find the best match between
    them\pause
  \item We can measure this distance in terms of an \emph{edit
      distance} or \emph{Levenshtein} distance which measures the
    minimum number of moves to go from one sequence to another\pause
  \item There are two cases that are commonly considered
    \begin{itemize}
    \item Global matching: where we match the two complete strings\pause
    \item Local matching: we find the sub-string with the best match (we
      treat this in the next lecture)\pause
    \end{itemize}
  \end{itemize}
\end{PauseHighLight}

\end{slide}


%%%%%%%%%%%%%%%%%%%%%%% Next Slide %%%%%%%%%%%%%%%%%%%%%%%
\Outline % Dynamic Programming
%%%%%%%%%%%%%%%%%%%%%%% Next Slide %%%%%%%%%%%%%%%%%%%%%%%

\begin{slide}
\section{Algorithms for Matching}

\begin{PauseHighLight}
  \begin{itemize}
  \item We need efficient algorithms for matching\pause
  \item The number of possible inexact matches is exponential in the size of the
    string so we can't exhaustively try all matches\pause
  \item Fortunately dynamical programming provides a relatively
    efficient algorithm\pause
  \item Dynamic programming is very important to a huge number of
    sequence-based algorithms\pause
  \item We first quickly remind you of dynamic programming\pause
  \end{itemize}
\end{PauseHighLight}

\end{slide}
%%%%%%%%%%%%%%%%%%%%%%% Next Slide %%%%%%%%%%%%%%%%%%%%%%%


\begin{slide}
\section{Dynamic Programming}

\begin{PauseHighLight}
  \begin{itemize}
  \item One of the most powerful strategies for solving optimisation
    problems is \emph{dynamic programming}\pause
    \begin{itemize}\squeeze
    \item Build a set of optimal partial solutions\pause
    \item Increase the size of the partial solutions until you have a
      full solution\pause
    \item Each step uses the set of optimal partial solutions found in
      the previous step\pause
    \end{itemize}
  \item Developed by Richard Bellman in the early 1950's\pause
  \item The name is unfortunate as it doesn't have much to do with
    programming\pause 
  \end{itemize}
\end{PauseHighLight}

\end{slide}

%%%%%%%%%%%%%%%%%%%%%%% Next Slide %%%%%%%%%%%%%%%%%%%%%%%

\begin{slide}
\section[-2]{A Toy Problem}

\begin{PauseHighLight}
  \begin{itemize}
  \item Consider the problem of find a minimum cost path from point
    (0,0) to (8,0) on the lattice\pause
    \begin{center}\color{TextColor}
      
\includegraphics[height=8cm]{lattice_prob.mps}
    \end{center}\color{TwoColor}
  \item The costs of traversing each link is shown in red\pause
  \item The cost of a path is the sum of weights on each link\pause
  \end{itemize}
\end{PauseHighLight}

\end{slide}

%%%%%%%%%%%%%%%%%%%%%%% Next Slide %%%%%%%%%%%%%%%%%%%%%%%

\begin{slide}
\section[-0.5]{Brute Force}

\begin{PauseHighLight}
  \begin{itemize}
  \item The obvious brute force strategy is to try every path\pause
  \item For a problem with $n$ steps we require $n/2$ to be diagonally
    up and $n/2$ to be diagonally down\pause
  \item The total number of paths is
    \begin{align*}
      \binom{n}{n/2} \approx \sqrt{\frac{2}{\pi\,n}}\; 2^n\pause
    \end{align*}
  \item For the problem shown above with $n=8$ there are 70
    paths\pause
  \item For a problem with $n=100$ there are $1.01\times10^{29}$
    paths\pause
  \end{itemize}
\end{PauseHighLight}

\end{slide}

%%%%%%%%%%%%%%%%%%%%%%% Next Slide %%%%%%%%%%%%%%%%%%%%%%%

\begin{slide}
\section[-2]{Building a Solution}

\begin{PauseHighLight}
  \begin{itemize}
  \item We can solve this problem efficiently using dynamic programming
    by considering optimal paths of shorter length\pause
  \item Let $c_{(i,j)}$ denote the cost of the optimal path to node
    $(i,j)$\pause
  \item We denote the weights between two points on the lattice by
    $w_{(i,j)(i+1,j\pm1)}$
    \begin{center}
      \includegraphics[height=4cm]{lattice_prob_notation}\pause
    \end{center}
  \item We define the initial cost $c_{(0,0)}=0$\pause
  \end{itemize}
\end{PauseHighLight}

\end{slide}

%%%%%%%%%%%%%%%%%%%%%%% Next Slide %%%%%%%%%%%%%%%%%%%%%%%

\begin{slide}
\section[-0.5]{Forward Algorithm}

\begin{PauseHighLight}
  \begin{itemize}
  \item Suppose we know the optimal costs for all the edge in column
    $i$\pause
  \item Our task is to find the optimal cost at column $i+1$\pause
  \item If we consider the sites in the lattice then the optimal cost
    will be
    \begin{align*}
      c_{(i+1,j)} = \min\!\left( \strut c_{(i,j+1)}+w_{(i,j+1)(i+1,j)},\,\,
        c_{(i,j-1)}+w_{(i,j-1)(i+1,j)}\right)\pause
    \end{align*}
  \item This is the defining equation in dynamic programming\pause
  \item We have to treat the boundary sites specially, but this is just
    book-keeping\pause
  \end{itemize}
\end{PauseHighLight}

\end{slide}

%%%%%%%%%%%%%%%%%%%%%%% Next Slide %%%%%%%%%%%%%%%%%%%%%%%

\begin{slide}
\section[-2]{Example}
\pb
\begin{center}
  \multiinclude[graphics={height=15cm}]{dynamicprog}\pause
\end{center}
\end{slide}

%%%%%%%%%%%%%%%%%%%%%%% Next Slide %%%%%%%%%%%%%%%%%%%%%%%

\begin{slide}
\section{Backward Algorithm}

\begin{PauseHighLight}
  \begin{itemize}
  \item Having found the optimal costs $c_{(i,j)}$ we can find the
    optimal path starting from $(n,0)$\pause
  \item At each step we have a choice of going up or down\pause
  \item We choose the direction which satisfies the constraint
    \begin{align*}
      c_{(i,j)} = c_{(i-1,j\pm1)}+w_{(i-1,j\pm1)(i,j)}\pause
    \end{align*}
  \item If both directions satisfy the constraint we have more than one
    optimal path\pause
  \end{itemize}
\end{PauseHighLight}

\end{slide}

%%%%%%%%%%%%%%%%%%%%%%% Next Slide %%%%%%%%%%%%%%%%%%%%%%%

\begin{slide}
\section{Time Complexity}

\begin{PauseHighLight}
  \begin{itemize}
  \item In our dynamic programming solution we had to compute the cost
    $c_{(i,j)}$ at each lattice point\pause
  \item There were $(n/2+1)^2$ lattice point\pause
  \item It took constant time to compute each cost so the total time to
    perform the forward algorithm was $\Theta(n^2)$\pause
  \item The time complexity of the backward algorithm was
    $\Theta(n)$\pause
  \item This compares with $\exp(\Theta(n))$ for the brute force
    algorithm\pause 
  \end{itemize}
\end{PauseHighLight}

\end{slide}

%%%%%%%%%%%%%%%%%%%%%%% Next Slide %%%%%%%%%%%%%%%%%%%%%%%
\Outline % Global Matching
%%%%%%%%%%%%%%%%%%%%%%% Next Slide %%%%%%%%%%%%%%%%%%%%%%%

\begin{slide}
\section[-2]{Global Matching}

\begin{PauseHighLight}
  \begin{itemize}
  \item We return to the problem of global matching\pause
  \item In biology this is often known as the Needleman-Wunsch problem
    after the author that proposed the problem and a (not very
    efficient) solution\pause
  \item We denote the cost of a substitution, $a\rightarrow b$, by
    $s(a,b)$\pause
  \item For proteins we can take this to be our Blosum50 matrix\pause
  \item There are different substitution matrices e.g. Blosum60 used for
    closely related species\pause
  \item We will assume the cost of an insertion or deletion is -8\pause
  \end{itemize}
\end{PauseHighLight}

\end{slide}

%%%%%%%%%%%%%%%%%%%%%%% Next Slide %%%%%%%%%%%%%%%%%%%%%%%

\begin{slide}
\section[-2]{Problem Definition}

\begin{PauseHighLight}
  \begin{itemize}
  \item We are given two strings
    \begin{align*}
      \bm{a} &= a_1\,a_2\,\ldots\,a_m &
      \bm{b} &= b_1\,b_2\,\ldots\,b_n\pause
    \end{align*}
  \item We want to find the minimum edit distance measuring how to
    change $\bm{a}$ to $\bm{b}$ where the price of a substitution is
    $s(a,b)$ and of an insertion or deletion $d=-8$\pause
  \item To compute this we compute a tableau showing the best possible
    moves using the dynamic programming formula\pause
    \begin{align*}
      F(i,j) = \max \left\{
      \begin{matrix}
        F(i-1,j-1) + s(a_i, b_j)\\
        F(i-1,j) + d\\
        F(i,j-1) + d
      \end{matrix}\right.\pause
    \end{align*}
  \end{itemize}
\end{PauseHighLight}

\end{slide}

%%%%%%%%%%%%%%%%%%%%%%% Next Slide %%%%%%%%%%%%%%%%%%%%%%%

\begin{slide}
\section[-2]{Global Alignment}
\pause
\pb
\pauselevel{=1}
\begin{center}
  \multipdf[width=\linewidth]{needlemanWunsch}\pause
\end{center}
\end{slide}

%%%%%%%%%%%%%%%%%%%%%%% Next Slide %%%%%%%%%%%%%%%%%%%%%%%

\begin{slide}
\section{Space Complexity}

\begin{PauseHighLight}
  \begin{itemize}
  \item For two strings of length $m$ and $n$ the time complexity of
    dynamic programming is the size of the tableau $O(m\,n)$\pause
  \item This is massively faster than exhaustive search\pause
  \item Unfortunately in biology we often want to match two very long
    strings (particularly when looking at DNA)\pause
  \item The space complexity is then more of an issue than time
    complexity (we don't have enough memory to store the tableau)\pause
  \item Algorithms exist that use linear space complexity\pause
  \item Can still be too slow for DNA\pause
  \end{itemize}
\end{PauseHighLight}

\end{slide}


%%%%%%%%%%%%%%%%%%%%%%% Next Slide %%%%%%%%%%%%%%%%%%%%%%%

\begin{slide}
\section{Summary}

\begin{PauseHighLight}
  \begin{itemize}
  \item Comparing sequences is done routinely to understand function and
    evolution\pause
  \item The dynamic programming strategy is widely used for inexact matching\pause
  \item For global matching the dynamic programming solution is really
    quite easy\pause
  \item We will also see that it is readily extendable\pause
  \end{itemize}
\end{PauseHighLight}


\end{slide}
